\chapter{完全失相干的独立事件极限}
\label{app:incoherent-event-chapter}

\label{app:incoherent-event-model}

本附录讨论一个非常具体的统计极限：把吸收与发射看成彼此独立、相位完全丢失的单光子事件。它不是任意开放系统退相干的普适主方程，只是与相干 Bessel 边带作对照的独立事件模型。其输入仍然是正文同一个单光子弱耦合尺度，Maxwell 与轨迹投影部分完全不变。


正文的 $P_n=J_n^2(2|\beta|)$ 来自相干相位调制：各阶吸收/发射振幅在取模平方之前保持相位关系。为比较完全失相干的极限，另设一个独立事件模型：每次事件只交换一个光子，吸收与发射相互独立且相位信息完全丢失。弱耦合式~\eqref{eq:weak_single} 给出总事件平均数
\begin{equation}
\boxed{
\bar N=P(+\hbar\omega)+P(-\hbar\omega)
=2|\beta|^2
=2\left(\frac{e|\Ftilde|}{2\hbar\omega}\right)^2.}
\label{eq:Nbar-incoherent}
\end{equation}
为避免与电子质量 $m$、球多极 $\ell$ 冲突，记吸收、发射事件数分别为 $N_+$、$N_-$。设
\begin{equation}
N_+\sim\operatorname{Poisson}(\bar N/2),
\qquad
N_-\sim\operatorname{Poisson}(\bar N/2),
\end{equation}
且二者独立。净边带阶数为
\begin{equation}
n=N_+-N_-.
\end{equation}
直接卷积两个 Poisson 分布。对 $n\ge0$，令 $N_-=j$、$N_+=n+j$，则
\begin{align}
P_n^{\rm incoh}(\bar N)
&=\sum_{j=0}^{\infty}
\ee^{-\bar N/2}\frac{(\bar N/2)^{n+j}}{(n+j)!}
\ee^{-\bar N/2}\frac{(\bar N/2)^j}{j!}\\
&=\ee^{-\bar N}
\sum_{j=0}^{\infty}
\frac{1}{j!(n+j)!}
\left(\frac{\bar N}{2}\right)^{n+2j}.
\end{align}
利用修正 Bessel 函数的级数定义
\begin{equation}
I_n(\bar N)=\sum_{j=0}^{\infty}
\frac{1}{j!(n+j)!}
\left(\frac{\bar N}{2}\right)^{n+2j},
\end{equation}
以及 $I_{-n}=I_n$，得到 Skellam 分布
\begin{equation}
\boxed{
P_n^{\rm incoh}(\bar N)=\ee^{-\bar N}I_{|n|}(\bar N).}
\label{eq:Skellam}
\end{equation}
这种构造与“总事件数服从 Poisson，再以 $1/2$ 的概率二项拆成吸收/发射”数学等价，但当前写法避免再引入会与质量或多极阶数冲突的计数符号。

对于脉冲，电子密度为
\begin{equation}
\rho(z')=\frac{1}{\sqrt{2\pi}v_0\sigma_{\mathrm e}}
\exp\!\left[-\frac{z'^2}{2v_0^2\sigma_{\mathrm e}^2}\right].
\end{equation}
电子包络中位置 \(z'\) 的那一小段在
\begin{equation}
t=-\frac{z'}{v_0}
\end{equation}
到达位于 \(z=0\) 的纳米结构。若激光强度包络为
\begin{equation}
I(0,t)=I_0\exp\!\left[-\frac{(t-\tau)^2}{2\sigma_{\mathrm L}^2}\right],
\end{equation}
则该电子片段看到的强度为
\begin{equation}
I\!\left(0,-\frac{z'}{v_0}\right)
=I_0\exp\!\left[-\frac{(z'+v_0\tau)^2}{2v_0^2\sigma_{\mathrm L}^2}\right].
\end{equation}
因为 \(\bar N\propto|\Ftilde|^2\propto I\)，故
\begin{equation}
\boxed{
\bar N(z')=\bar N_0
\exp\!\left[-\frac{(z'+v_0\tau)^2}{2v_0^2\sigma_{\mathrm L}^2}\right].}
\label{eq:Nbar-z}
\end{equation}
总的第 \(n\) 阶比例为
\begin{equation}
f_n=\int_{-\infty}^{\infty}\rho(z')
\ee^{-\bar N(z')}I_n(\bar N(z'))\dd z'.
\label{eq:fn_incoherent}
\end{equation}
为了得到显式双重级数，同时展开
\begin{equation}
\ee^{-\bar N}=\sum_{j=0}^{\infty}\frac{(-\bar N)^j}{j!},
\qquad
I_n(\bar N)=\sum_{k=0}^{\infty}
\frac{1}{k!(n+k)!}\left(\frac{\bar N}{2}\right)^{n+2k}.
\end{equation}
于是
\begin{align}
f_n
={}&\sum_{j,k\ge0}
\frac{(-1)^j}{j!k!(n+k)!2^{n+2k}}
\bar N_0^{n+j+2k}\\
&\times\int_{-\infty}^{\infty}
\frac{\dd z'}{\sqrt{2\pi}v_0\sigma_{\mathrm e}}
\exp\!\left[-\frac{z'^2}{2v_0^2\sigma_{\mathrm e}^2}
-\frac{(n+j+2k)(z'+v_0\tau)^2}
{2v_0^2\sigma_{\mathrm L}^2}\right].
\label{eq:incoherent-series-coefficients}
\end{align}
令
\begin{equation}
s=n+j+2k.
\end{equation}
这仍是两个 Gaussian 的乘积，但这里把积分完整算出。记
\begin{equation}
A\equiv\frac{1}{\sigma_{\mathrm e}^2}+\frac{s}{\sigma_{\mathrm L}^2}
=\frac{\sigma_{\mathrm L}^2+s\sigma_{\mathrm e}^2}
{\sigma_{\mathrm e}^2\sigma_{\mathrm L}^2}.
\end{equation}
指数可写成
\begin{align}
&-\frac{z'^2}{2v_0^2\sigma_{\mathrm e}^2}
-\frac{s(z'+v_0\tau)^2}{2v_0^2\sigma_{\mathrm L}^2}\\
&=-\frac{1}{2v_0^2}
\left[
Az'^2+\frac{2sv_0\tau}{\sigma_{\mathrm L}^2}z'
+\frac{sv_0^2\tau^2}{\sigma_{\mathrm L}^2}
\right]\\
&=-\frac{A}{2v_0^2}
\left(z'+\frac{sv_0\tau}{A\sigma_{\mathrm L}^2}\right)^2
-\frac{s\tau^2}{2(\sigma_{\mathrm L}^2+s\sigma_{\mathrm e}^2)}.
\end{align}
因此
\begin{align}
&\int_{-\infty}^{\infty}
\frac{\dd z'}{\sqrt{2\pi}v_0\sigma_{\mathrm e}}
\exp\!\left[-\frac{z'^2}{2v_0^2\sigma_{\mathrm e}^2}
-\frac{s(z'+v_0\tau)^2}{2v_0^2\sigma_{\mathrm L}^2}\right]\\
&=\frac{1}{\sqrt{2\pi}v_0\sigma_{\mathrm e}}
\sqrt{\frac{2\pi v_0^2}{A}}
\exp\!\left[-\frac{s\tau^2}{2(\sigma_{\mathrm L}^2+s\sigma_{\mathrm e}^2)}\right]\\
&=\frac{1}{\sqrt{1+s(\sigma_{\mathrm e}/\sigma_{\mathrm L})^2}}
\exp\!\left[
-\frac{s(\tau/\sigma_{\mathrm L})^2}
{2+2s(\sigma_{\mathrm e}/\sigma_{\mathrm L})^2}
\right].
\label{eq:incoherent-gaussian-integral-explicit}
\end{align}
把系数整理为
\begin{equation}
\boxed{
c_{jk}^n
=\frac{(-2)^j}{j!k!(n+k)!}
\left(\frac{\bar N_0}{2}\right)^{n+j+2k}}
\label{eq:cnjk_incoherent}
\end{equation}
后得到
\begin{equation}
\boxed{
f_n
=\sum_{j,k\ge0}c_{jk}^n
\frac{1}{\sqrt{1+s(\sigma_{\mathrm e}/\sigma_{\mathrm L})^2}}
\exp\!\left[
-\frac{s(\tau/\sigma_{\mathrm L})^2}
{2+2s(\sigma_{\mathrm e}/\sigma_{\mathrm L})^2}
\right],
\qquad s=n+j+2k.}
\label{eq:incoherent-delay-series}
\end{equation}
再对 \((b,\varphi)\) 做与式~\eqref{eq:Dfactor_start} 完全相同的空间平均，只需把相干模型中的
\(
s=n+j+k
\)
替换成
\(
s=n+j+2k
\)。于是球形几何
\begin{equation}
\boxed{
d_{jk}^{n,(sph)}=
\frac{\delta\left[(2sa+\delta)-
\ee^{-2s(w-a)/\delta}(2sw+\delta)\right]}
{2\sqrt\pi\,s^2(w^2-a^2)}
\frac{\Gamma(s+\frac12)}{\Gamma(s+1)},
\qquad s=n+j+2k,}
\label{eq:d_incoherent_sphere}
\end{equation}
圆柱几何
\begin{equation}
\boxed{
d_{jk}^{n,(cyl)}=
\frac{\delta[1-\ee^{-2s(w-a)/\delta}]}
{2s(w-a)},
\qquad s=n+j+2k.}
\label{eq:d_incoherent_cyl}
\end{equation}
当 $s=0$（只可能发生在 $n=j=k=0$）时，两种闭式都按连续极限取 $d_{00}^{0}=1$；不能把含 $s^{-1}$ 或 $s^{-2}$ 的表达式机械代入 $s=0$。最终
\begin{equation}
\boxed{
P_{\rm EEGS}^{\rm incoh}(n)
=\sum_{j,k\ge0}d_{jk}^n c_{jk}^n
\frac{1}{\sqrt{1+s(\sigma_{\mathrm e}/\sigma_{\mathrm L})^2}}
\exp\!\left[
-\frac{s(\tau/\sigma_{\mathrm L})^2}
{2+2s(\sigma_{\mathrm e}/\sigma_{\mathrm L})^2}
\right].}
\label{eq:incoherent-observed-spectrum}
\end{equation}
由此得到完整的非相干时空平均概率。与相干 STS 的关键差别是：相干理论的幅度先求和再取模平方，因此出现 \(C_j^nC_k^{n*}\) 的干涉交叉项以及普通 Bessel 函数 \(J_n\) 的振荡；非相干模型直接对事件概率求和，得到 \(\ee^{-\bar N}I_n(\bar N)\)，不保留光学相位。


\end{document}
